Ten conditions were run with three seeds each at $B = 2{,}400$, and every final student was
evaluated greedily on the 1,500 test problems (Table~\ref{ca:tab:h3}).

\begin{table}[htbp]
\centering
\footnotesize
\setlength{\tabcolsep}{2.5pt}
\begin{tabular}{@{}lrccccc@{}}
\toprule
 & & \multicolumn{3}{c}{\textbf{Test accuracy}} & & \\
\cmidrule(l){3-5}
\textbf{Condition} & \textbf{Rounds} & \textbf{Countdown} & \textbf{MATH} & \textbf{Code} & \textbf{Utility [95\% CI]} & \textbf{Diff. [95\% CI]} \\
\midrule
Untrained student & -- & 0.218 & 0.434 & 0.336 & 0.329 & \\
\midrule
Qwen2.5-3B only & 400 & 0.555 & 0.418 & 0.352 & 0.442 [0.416, 0.468] & $-$0.052 [$-$0.083, $-$0.024] \\
Phi-4-mini only & 400 & 0.630 & 0.438 & 0.352 & 0.473 [0.440, 0.508] & $-$0.020 [$-$0.049, +0.009] \\
Qwen2.5-7B only & 364 & 0.698 & 0.423 & 0.360 & 0.494 [0.468, 0.519] & -- \\
Qwen-Coder-7B only & 360 & 0.656 & 0.421 & 0.358 & 0.478 [0.450, 0.507] & $-$0.015 [$-$0.032, +0.000] \\
Phi-4 only & 110 & 0.417 & 0.429 & 0.335 & 0.394 [0.370, 0.417] & $-$0.100 [$-$0.127, $-$0.073] \\
\midrule
Uniform & 223 & 0.556 & 0.437 & 0.343 & 0.445 [0.420, 0.471] & $-$0.048 [$-$0.068, $-$0.026] \\
Headroom, learned eff. & 368 & 0.711 & 0.435 & 0.355 & 0.500 [0.478, 0.523] & +0.007 [$-$0.011, +0.026] \\
Cost-aware Thompson & 400 & 0.508 & 0.435 & 0.346 & 0.430 [0.329, 0.505] & $-$0.064 [$-$0.165, +0.004] \\
\midrule
Correctness, no provider & 200 & 0.611 & 0.425 & 0.352 & 0.462 [0.437, 0.489] & $-$0.031 [$-$0.052, $-$0.011] \\
SFT on responses & -- & 0.373 & 0.421 & 0.322 & 0.372 [0.350, 0.393] & $-$0.122 [$-$0.144, $-$0.099] \\
\bottomrule
\end{tabular}
\caption{Test results of the main runs at budget $B = 2{,}400$, means over three seeds. Utility is the mean of the three domain accuracies. Intervals are 95\% bootstrap intervals over seeds and test items; the last column is the paired difference from the best single provider, Qwen2.5-7B. Rounds are gradient steps, capped at 400; SFT takes no rounds. Runs that reached the cap stopped before spending the whole budget; mean spend was lowest for Qwen2.5-3B alone (1,374 units).}
\label{ca:tab:h3}
\end{table}


\paragraph{The proposed selector matches the best single provider but does not beat it.}
The best single provider is Qwen2.5-7B, with test utility 0.494. The headroom selector with
learned efficiency reaches 0.500, the highest of any condition and the most consistent
across seeds (0.503, 0.498 and 0.500), but its paired advantage of $+0.007$ has a 95\%
interval from $-0.011$ to $+0.026$. H3 is therefore not supported at this budget. The
selector did learn what the pilot version could not: it spent 38 to 58\% of its budget on
Countdown, where learning happens, and only 6 to 9\% on Phi-4. It did not reliably avoid the
degraded provider, which received 16 to 17\% of the budget in two seeds and none in the third.

\paragraph{Other conditions.} Cost-aware Thompson sampling reached 0.465 and 0.502 in two
seeds but collapsed in the third (0.322), where the student's Countdown answers shrank from
445 tokens to 38. Uniform allocation spent 59 to 66\% of its budget on Phi-4 and its degraded
copy and did no better than the cheapest single provider (0.445 against 0.442 for
Qwen2.5-3B). Phi-4 alone was the worst single provider,
because its price bought only 110 rounds. The correctness control, which uses ground-truth
labels and no provider for 200 rounds, reached 0.462, and SFT on provider responses was the
weakest condition at 0.372, below every on-policy run. As in the pilot, all of the gains are
on Countdown: no condition moved MATH or code test accuracy by more than 0.024 from the
untrained student.

These comparisons take the best static allocation over the single-provider runs only. The
planned per-domain allocation, which gives each domain its best provider, is not yet built
(Section~\ref{ca:sec:next}).
